\SourceLeaf{348}

\begin{center}
[III.]

\textbf{ZUR THEORIE DER KEGELSCHNITTE.}
\end{center}

Zunächst beziehen sich die Untersuchungen auf die Ellipse, sie lassen sich
aber mit geringfügigen Änderungen auf die Hyperbel übertragen.

\Needspace{8\baselineskip}\begin{center}
[1.]
\end{center}

Es seien $m$, $n$ zunächst zwei reelle Zahlen, $\rho$ eine Größe von der
Form $\cos\varphi + i\sin\varphi$, sodass $\rho$ alle seine Werte
durchläuft, während $\varphi$ von $0^0$ bis $360^0$ wächst.

Es sei ferner
\[
r = m\rho + \dfrac{n}{\rho},
\]
und $R$ der Punkt, der die komplexe Größe $r$ darstellt; $R$ ist also
ein unbestimmter Punkt auf einer Ellipse, deren Hauptachsen $m+n$, $m-n$ auf
der reellen beziehungsweise imaginären Achse liegen, während ihr
Mittelpunkt mit dem Nullpunkt zusammenfällt.

Den vier Werten von $\rho$, nämlich $\alpha$, $\beta$, $\gamma$, $\delta$,
entsprechen die vier Werte von $r$, nämlich $a$, $b$, $c$, $d$, oder die vier
Punkte $A$, $B$, $C$, $D$.

\Needspace{7\baselineskip}Schließlich setze man
\[
\dfrac{-\alpha\beta\gamma m^3 + (\alpha+\beta+\gamma) m m n
  - \delta(\alpha\beta+\alpha\gamma+\beta\gamma) m n n + \delta n^3}
  {(1-\alpha\beta\gamma\delta) m n} = t
\]
und nenne den Punkt, auf den sich die komplexe Zahl $t$ bezieht, $T$.

Man sieht leicht, dass, wenn von den vier Punkten $A$, $B$, $C$, $D$ drei
fest sind und der vierte den Kegelschnitt durchläuft, $T$ eine Gerade
durchläuft.

\Needspace{9\baselineskip}Eine einfache Rechnung ergibt
\[
t - a = -\dfrac{(\alpha\beta m - n)(\alpha\gamma m - n)(m - \alpha\delta n)}
  {\alpha(1-\alpha\beta\gamma\delta) m n}
\]
\[
= -\dfrac{(a-b)(a-c)}{a-d} \cdot
  \dfrac{\alpha\beta\gamma(\alpha-\delta)\left(mm -
  \left(\alpha\delta+\dfrac{1}{\alpha\delta}\right) m n + n n\right)}
  {(\alpha-\beta)(\alpha-\gamma)(1-\alpha\beta\gamma\delta) m n}.
\]
